\documentclass[a4paper]{book}

\setlength{\headheight}{14.5pt}

% --- 1. ENCODAGE ET LANGUE (À mettre en premier) ---
\usepackage[T1]{fontenc}
\usepackage[utf8]{inputenc}
\usepackage[french]{babel}

% --- 2. POLICES ET SYMBOLES ---
% fourier : police élégante pour les maths (lmodern supprimé car doublon)
\usepackage{fourier}
\usepackage{amsfonts, amssymb, amsthm}
\usepackage[right]{eurosym}
\usepackage[np]{numprint}
\usepackage[nointegrals]{wasysym}

% --- 3. GÉOMÉTRIE ET MISE EN PAGE ---
\usepackage[left=1cm,right=1cm,top=1.5cm,bottom=1.5cm]{geometry}
\usepackage[dvipsnames,svgnames]{xcolor}
\usepackage[most,skins,breakable]{tcolorbox}
\usepackage{tagging}

% --- 4. GRAPHISMES ET DIAGRAMMES ---
\usepackage{graphicx}
\usepackage{pgfplots}
\pgfplotsset{compat=1.18}
\usetikzlibrary{babel, shapes, fit} % 'babel' évite les bugs avec TikZ

% --- 5. TABLEAUX ---
\usepackage{array, booktabs, multirow, tabularx, longtable}

% --- 6. STRUCTURE ET TITRES ---
% titlesec est plus moderne que sectsty, on garde uniquement celui-ci
\usepackage{titlesec}
\usepackage{enumitem}
\usepackage[normalem]{ulem} % 'normalem' évite que \emph devienne souligné

% --- 7. HYPERLIENS (toujours en dernier ou presque) ---
\usepackage{hyperref}

\setlength{\parindent}{0pt}
\frenchbsetup{StandardItemLabels=true}

% Le package ultime
\usepackage{ProfCollege}


% =================================================================
% COMMANDES GÉNÉRALES
% =================================================================

% Une ligne de points (usage : \points[15])
\newcommand{\points}[1][1]{%
    \foreach \n in {1,...,#1} { \ldots}%
}

% Plusieurs lignes de points (usage : \lignesPoints{3})
\newcommand{\lignesPoints}[1]{%
    \foreach \n in {1,...,#1} {%
        \dotfill \newline
    }%
}


% =================================================================
% ZONES RÉPONSE
% =================================================================

% Zone réponse universelle (usage : \begin{ZoneReponse}{4cm})
\newtcolorbox{ZoneReponse}[1]{
    enhanced,
    colback=white,
    colframe=gray!50,
    arc=2mm,
    boxrule=0.5mm,
    width=\textwidth,
    height=#1,
    valign=top,
    fonttitle=\bfseries,
    colbacktitle=gray!20,
    coltitle=black,
    attach title to upper,
    after skip=10pt plus 2pt minus 2pt
}


% =================================================================
% COMMANDES MATHÉMATIQUES COLLÈGE
% =================================================================

% Ensembles de nombres
\newcommand{\N}{\ensuremath{\mathbf{N}}}
\newcommand{\Z}{\ensuremath{\mathbf{Z}}}
\newcommand{\D}{\ensuremath{\mathbf{D}}}
\newcommand{\Q}{\ensuremath{\mathbf{Q}}}
\newcommand{\R}{\ensuremath{\mathbf{R}}}

% Valeur absolue
\newcommand{\abs}[1]{\left\lvert #1 \right\rvert}

% Arc géométrique (usage : $\arc{AB}$)
% Note : utilise la fonte tipa pour le symbole arc
\makeatletter
\DeclareFontFamily{U}{tipa}{}
\DeclareFontShape{U}{tipa}{m}{n}{<->tipa10}{}
\newcommand{\arc@char}{{\usefont{U}{tipa}{m}{n}\symbol{62}}}
\newcommand{\arc}[1]{\mathpalette\arc@arc{#1}}
\newcommand{\arc@arc}[2]{%
  \sbox0{$\m@th#1#2$}%
  \vbox{\hbox{\resizebox{\wd0}{\height}{\arc@char}}\nointerlineskip\box0}%
}
\makeatother

% Comparaison : boîte vide entre deux nombres pour écrire le symbole
% Usage : \compare{102}{120}  ou  \compare[0.5]{\np{132999999}}{\np{133000000}}
% Argument optionnel : ignoré (compatibilité ascendante), la boîte s'adapte
\newcommand{\compare}[3][0]{%
    #2 \kern0.4em
    \tikz[baseline=-0.6ex]
        \node[rectangle, draw=gray!60, thick, dotted, minimum height=0.6cm,
              minimum width=0.6cm, inner sep=2pt] {};%
    \kern0.4em #3%
}

% Encadrement : deux boîtes vides autour d'un nombre
% Usage : \encadre{\np{1389}}  ou  \encadre[1.8]{\np{652176}}
% Argument optionnel : largeur des boîtes en cm (défaut 1.5)
\newcommand{\encadre}[2][1.5]{%
    \tikz[baseline=-0.6ex]
        \node[rectangle, draw=gray!60, thick, dotted, minimum height=0.6cm,
              minimum width=#1cm, inner sep=2pt] {};%
    \kern0.3em $<$ #2 $<$ \kern0.3em
    \tikz[baseline=-0.6ex]
        \node[rectangle, draw=gray!60, thick, dotted, minimum height=0.6cm,
              minimum width=#1cm, inner sep=2pt] {};%
}

% Diagramme de proportionnalité simple
% Usage : \diagrammesimple{A}{B}{×3}
\newcommand{\diagrammesimple}[3]{%
\begin{tikzpicture}[baseline=(A.base)]
    \node(A) at (0,0) {#1};
    \node(B) at (5,0) {#2};
    \draw[->, dashed, line width=0.8mm, green!50!black]
        (A) to[bend left=20] node[midway, above=2pt] {#3} (B);
\end{tikzpicture}}

% Diagramme avec flèche aller-retour
% Usage : \diagrammesimplereciproque{A}{B}{×3}{÷3}
\newcommand{\diagrammesimplereciproque}[4]{%
\begin{tikzpicture}[baseline=(A.base)]
    \node(A) at (0,0) {#1};
    \node(B) at (5,0) {#2};
    \draw[->, dashed, line width=0.8mm, green!50!black]
        (A) to[bend left=20] node[midway, above=2pt] {#3} (B);
    \draw[->, dashed, line width=0.8mm, blue!50!black]
        (B) to[bend left=20] node[midway, below=2pt] {#4} (A);
\end{tikzpicture}}
\newcommand{\maClasse}{3e}

%% --- PACKAGES PROPRES AUX COURS ---
\usepackage{fancyhdr}
\usepackage{lastpage}
\usepackage[column=X]{cellspace}

%% --- VARIABLES ---
\newcommand{\monTitre}{Titre du cours} % À redéfinir dans chaque fichier : \renewcommand{\monTitre}{Mon titre}

%% --- FORMAT DES TITRES (titlesec) ---

\titleformat{\chapter}[frame]
{\color{red!85!white}\normalsize}
{\filright\sffamily\Large\enspace Chapitre \thechapter\enspace}
{8pt}
{\color{red!85!white}\rule{0pt}{30pt}\sffamily\Huge\bfseries\filcenter}
\titlespacing{\chapter}{0pc}{0pc}{2pc}

\titleformat{\section}
{\color{green!55!black}\normalfont\Large\bfseries}
{\thesection}
{1em}
{}

% 3. Subsections (Cyan)
% J'ai retiré le \hspace{1em} qui poussait le texte
\titleformat{\subsection}[block]
{\color{cyan!65!black}\normalfont\large\bfseries}
{\thesubsection}
{1em}
{}
% J'ai mis le premier paramètre à 0pt pour supprimer l'indentation de bloc
\titlespacing{\subsection}{0pt}{1.5ex}{1ex}

% 4. Subsubsections (Violet)
\titleformat{\subsubsection}[block]
{\color{violet}\normalfont\normalsize\bfseries}
{\thesubsubsection}
{1em}
{}
\titlespacing{\subsubsection}{0pt}{1ex}{1ex}

%% --- NUMÉROTATION ---
\renewcommand{\thesection}{\Roman{section}.}
\renewcommand{\thesubsection}{\arabic{subsection})}
\renewcommand{\thesubsubsection}{\alph{subsubsection})}
\setcounter{secnumdepth}{3}

%% --- EN-TÊTES ET PIEDS DE PAGE ---
\pagestyle{fancy}
\renewcommand\headrulewidth{1pt}
\fancyhead[LO,RE]{\maClasse}
\fancyhead[RO,LE]{\monTitre}
\renewcommand\footrulewidth{1pt}
\fancyfoot[LO,RE]{Année 2025-2026}
\fancyfoot[RO,LE]{\textbf{Page \thepage/\pageref{LastPage}}}

%% --- RÉGLAGES PARAGRAPHES ---
\setlength{\cellspacebottomlimit}{4pt}
\setlength{\cellspacetoplimit}{4pt}
\setlength{\parindent}{2pt}
\setlength{\tabcolsep}{0.5cm}

%% --- COMPTEURS ET BOITES TCOLORBOX ---

%% --- COMPTEURS AVEC RÉINITIALISATION À CHAQUE CHAPITRE ---

% On ajoute [chapter] à la fin pour que le compteur reparte à 0 à chaque \chapter
\newcounter{dfn}[chapter]
\newcounter{thm}[chapter]
\newcounter{cvn}[chapter]
\newcounter{prp}[chapter]
\newcounter{ex}[chapter]
\newcounter{mtd}[chapter]
\newcounter{rq}[chapter]
\newcounter{demo}[chapter]
\newcounter{appli}[chapter]
\newcounter{exo}[chapter]
\newcounter{lem}[chapter]

%% --- BLOC COMPLET DES BOITES TCOLORBOX ---

% --- DÉFINITIONS (Bleu) ---
\newtcolorbox{definition}{breakable,colback=white!6!white,colframe=blue!85!white,fonttitle=\bfseries,title=Définition \thedfn, code={\stepcounter{dfn}}, arc=3mm}
\newtcolorbox{definitionTitre}[1]{breakable,colback=white!6!white,colframe=blue!85!white,fonttitle=\bfseries,title=Définition \thedfn\ - #1, code={\stepcounter{dfn}}, arc=3mm}
\newtcolorbox{definitions}{breakable,colback=white!6!white,colframe=blue!85!white,fonttitle=\bfseries,title=Définitions \thedfn, code={\stepcounter{dfn}}, arc=3mm}

% --- PROPRIÉTÉS & THÉORÈMES (Rouge/Orange) ---
\newtcolorbox{theoreme}{breakable,colback=red!6!white,colframe=red!85!white,fonttitle=\bfseries,title=Théorème \thethm, code={\stepcounter{thm}}, arc=3mm}
\newtcolorbox{theoremeTitre}[1]{breakable,colback=red!6!white,colframe=red!85!white,fonttitle=\bfseries,title=Théorème \thethm\ - #1, code={\stepcounter{thm}}, arc=3mm}
\newtcolorbox{propriete}{breakable,colback=red!6!white,colframe=red!85!white,fonttitle=\bfseries,title=Propriété \theprp, code={\stepcounter{prp}}, arc=3mm}
\newtcolorbox{proprieteTitre}[1]{breakable,colback=red!6!white,colframe=red!85!white,fonttitle=\bfseries,title=Propriété \theprp\ - #1, code={\stepcounter{prp}}, arc=3mm}
\newtcolorbox{proprietes}{breakable,colback=red!6!white,colframe=red!85!white,fonttitle=\bfseries,title=Propriétés \theprp, code={\stepcounter{prp}}, arc=3mm}
\newtcolorbox{lemme}{breakable,colback=red!6!white,colframe=orange!85!white,fonttitle=\bfseries,title=Lemme \theprp, code={\stepcounter{prp}}, arc=3mm}
\newtcolorbox{convention}{breakable,colback=orange!6!white,colframe=orange!85!white,fonttitle=\bfseries,title=Convention \thecvn, code={\stepcounter{cvn}}, arc=3mm}

% --- EXEMPLES (Vert) — argument optionnel pour le titre ---
% Usage : \begin{exemple}             -> "Exemple 1"
%         \begin{exemple}[Mon titre]  -> "Exemple 1 - Mon titre"
\newtcolorbox{exemple}[1][]{
  enhanced,
  breakable,
  nobeforeafter,
  % --- ON FORCE TOUT À ZÉRO ---
  boxrule=0pt,        % Épaisseur du cadre global à 0
  titlerule=0pt,      % Supprime la ligne sous le titre
  sharp corners,      % Important pour les coupures propres
  % ----------------------------
  colback=white,
  colframe=white,     % On met le cadre en blanc (invisible) au cas où
  frame empty,        % On cache le cadre
  %
  borderline west={0.5mm}{0mm}{green!55!black}, % Seule la barre verte reste
  %
  % Gestion de la coupure
  toprule at break=0pt,
  bottomrule at break=0pt,
  %
  title={Exemple \theex\ifx&#1&\else\ - #1\fi},
  code={\stepcounter{ex}},
  attach boxed title to top left = {xshift = -3mm},
  boxed title style={size=small,colback=green!55!black,colframe=green!55!black, arc=3mm}
}

% --- REMARQUES (Orange) ---
\newtcolorbox{remarque}{enhanced, breakable, frame empty, nobeforeafter, borderline west={0.5mm}{0mm}{orange!85!white}, title=Remarque \therq, colback=white, code={\stepcounter{rq}}, attach boxed title to top left = {xshift = -3mm}, boxed title style={size=small,colback=orange!85!white,colframe=orange!85!white, arc=3mm}}
\newtcolorbox{remarques}{enhanced, breakable, frame empty, nobeforeafter, borderline west={0.5mm}{0mm}{orange!85!white}, title=Remarques \therq, colback=white, code={\stepcounter{rq}}, attach boxed title to top left = {xshift = -3mm}, boxed title style={size=small,colback=orange!85!white,colframe=orange!85!white, arc=3mm}}
\newtcolorbox{remarqueTitre}[1]{enhanced, breakable, frame empty, nobeforeafter, borderline west={0.5mm}{0mm}{orange!85!white}, title=Remarque \therq\ - #1, colback=white, code={\stepcounter{rq}}, attach boxed title to top left = {xshift = -3mm}, boxed title style={size=small,colback=orange!85!white,colframe=orange!85!white, arc=3mm}}

% --- MÉTHODES & APPLICATIONS (Bleu foncé) ---
\newtcolorbox{methode}{enhanced, breakable, frame empty, nobeforeafter, borderline west={0.5mm}{0mm}{blue!75!black}, title=Méthode \themtd, colback=white, code={\stepcounter{mtd}}, attach boxed title to top left = {xshift = -3mm}, boxed title style={size=small,colback=blue!75!black,colframe=blue!75!black, arc=3mm}}
\newtcolorbox{methodes}{enhanced, breakable, frame empty, nobeforeafter, borderline west={0.5mm}{0mm}{blue!75!black}, title=Méthodes \themtd, colback=white, code={\stepcounter{mtd}}, attach boxed title to top left = {xshift = -3mm}, boxed title style={size=small,colback=blue!75!black,colframe=blue!75!black, arc=3mm}}
\newtcolorbox{methodeTitre}[1]{enhanced, breakable, frame empty, nobeforeafter, borderline west={0.5mm}{0mm}{blue!75!black}, title=Méthode \themtd\ - #1, colback=white, code={\stepcounter{mtd}}, attach boxed title to top left = {xshift = -3mm}, boxed title style={size=small,colback=blue!75!black,colframe=blue!75!black, arc=3mm}}
\newtcolorbox{application}{enhanced, breakable, frame empty, nobeforeafter, borderline west={0.5mm}{0mm}{orange!75!black}, title=Application \themtd, colback=white, code={\stepcounter{mtd}}, attach boxed title to top left = {xshift = -3mm}, boxed title style={size=small,colback=orange!75!black,colframe=orange!75!black, arc=3mm}}

% --- DIVERS (Exercices, Python, Conséquences) ---
\newtcolorbox{exercice}{enhanced, breakable, frame empty, nobeforeafter, borderline west={0.5mm}{0mm}{green!55!black}, title=Exercice n°\theexo, colback=white, code={\stepcounter{exo}}, attach boxed title to top left = {xshift = -3mm}, boxed title style={size=small,colback=green!55!black,colframe=green!55!black, arc=3mm}}
\newtcolorbox{exerciceTitre}[1]{enhanced, breakable, frame empty, borderline west={0.5mm}{0mm}{green!55!black}, title=Exercice n°\theexo\ - #1, colback=white, code={\stepcounter{exo}}, attach boxed title to top left = {xshift = -3mm}, boxed title style={size=small,colback=green!55!black,colframe=green!55!black, arc=3mm}}
\newtcolorbox{consequences}{enhanced, breakable, frame empty, nobeforeafter, borderline west={0.5mm}{0mm}{ProcessBlue}, title=Conséquences, colback=white, attach boxed title to top left = {xshift = -3mm}, boxed title style={size=small,colback=ProcessBlue,colframe=ProcessBlue, arc=3mm}}
\newtcolorbox{programme}{enhanced, breakable, frame empty, nobeforeafter, borderline west={0.5mm}{0mm}{gray!55!black}, title=Programme Python, colback=white, attach boxed title to top left = {xshift = -3mm}, boxed title style={size=small,colback=gray!55!black,colframe=gray!55!black, arc=3mm}}


% Commande pour les démonstrations (incrémentation automatique)
\newcommand{\demonstration}{%
  \stepcounter{demo}%
  \tcbox[enhanced, 
         on line, 
         arc=2mm, 
         colback=gray!10!white, 
         colframe=gray!50!black, 
         size=small, 
         fontupper=\bfseries]
         {Démonstration \thedemo}%
  \enspace % Ajoute un petit espace après la boîte
}

%% --- FIN DU BLOC DES BOITES ---

% Activité (Numéro, Titre, Chapitre)
\newcommand{\enteteActivite}[3]{
    \noindent
    \begin{tabularx}{\textwidth}{@{}p{2cm} >{\centering\arraybackslash}X p{2cm}@{}}
        \textbf{Activité n°#1} & \textbf{#2} & \textbf{Chapitre #3} \\
    \end{tabularx}
}


% --- Commande pour les blocs mémo personnalisés ---
% Argument 1 : Couleur du cadre
% Argument 2 : Titre du bloc
% Argument 3 : Texte/Contenu
\newcommand{\blocMemo}[3]{%
    \begin{tcolorbox}[
        enhanced,
        breakable,
        title=#2,
        colback=#1!5,          % Fond très clair
        colframe=#1!75!black,  % Bordure
        fonttitle=\bfseries\sffamily,
        colbacktitle=#1!75!black,
        attach boxed title to top left={xshift=3mm, yshift=-3mm},
        boxed title style={size=small, arc=2mm, colframe=#1!75!black},
        arc=3mm,               % Bel arrondi
        boxrule=0.6mm,         % Trait un peu plus épais
        after skip=15pt,
        before skip=10pt
    ]
        \vspace{2mm} % Petit espace pour ne pas coller au titre boxé
        #3
    \end{tcolorbox}
}


% --- Nouvel environnement Culture Maths ---
\newtcolorbox{cultureMaths}{
    enhanced, 
    breakable, 
    frame empty, 
    nobeforeafter, 
    borderline west={0.5mm}{0mm}{gray!60!black}, % Bordure grise comme ton "programme"
    title=Le saviez-vous ?, 
    colback=white, 
    attach boxed title to top left = {xshift = -3mm}, 
    boxed title style={size=small,colback=gray!60!black,colframe=gray!60!black, arc=3mm},
    coltitle=white,
    fonttitle=\bfseries\sffamily,
    after skip=15pt,
    before skip=10pt
}

\begin{document}

\setcounter{chapter}{3}
\renewcommand{\monTitre}{Arithmétique}
\chapter{\monTitre}
\thispagestyle{fancy}


% ============================================================
\section{Multiples, diviseurs et nombres premiers}
% ============================================================

\subsection{Multiples et diviseurs}

    \begin{definition}
        Soit $a$ et $b$ deux entiers avec $b \neq 0$. On dit que :
        \begin{itemize}
            \item $b$ \textbf{divise} $a$ (ou que $b$ est un \textbf{diviseur} de $a$) si le reste de la division euclidienne de $a$ par $b$ est nul, c'est-à-dire si $a = b \times k$ pour un certain entier $k$.
            \item Dans ce cas, $a$ est un \textbf{multiple} de $b$.
        \end{itemize}
    \end{definition}

    \begin{exemple}
        \begin{enumerate}
            \item Lister tous les diviseurs de $24$ :
            \begin{ZoneReponse}{2cm}
            \end{ZoneReponse}

            \item $7$ est-il un diviseur de $\np{1001}$ ? Justifier.
            \begin{ZoneReponse}{2cm}
            \end{ZoneReponse}
        \end{enumerate}
    \end{exemple}

    \begin{remarque}
        Quelques critères de divisibilité utiles à connaître :
        \begin{itemize}
            \item Par $2$ : le chiffre des unités est $0, 2, 4, 6$ ou $8$.
            \item Par $3$ : la somme des chiffres est divisible par $3$.
            \item Par $5$ : le chiffre des unités est $0$ ou $5$.
            \item Par $9$ : la somme des chiffres est divisible par $9$.
            \item Par $10$ : le chiffre des unités est $0$.
        \end{itemize}
    \end{remarque}

    \begin{exemple}
        Parmi les nombres suivants, indiquer par quels entiers (parmi $2$, $3$, $5$, $9$, $10$) chacun est divisible :
        
        \medskip
        \noindent
        \begin{minipage}[t]{0.45\textwidth}
            \textbf{a)} $\np{2430}$
            \begin{ZoneReponse}{2cm}
            \end{ZoneReponse}
        \end{minipage}
        \hfill
        \begin{minipage}[t]{0.45\textwidth}
            \textbf{b)} $\np{5175}$
            \begin{ZoneReponse}{2cm}
            \end{ZoneReponse}
        \end{minipage}
    \end{exemple}


\subsection{Nombres premiers}

    \begin{definition}
        Un \textbf{nombre premier} est un entier qui admet \textbf{exactement deux diviseurs distincts} : $1$ et lui-même.
        
        \medskip
        Un entier supérieur à $1$ qui n'est pas premier est dit \textbf{composé}.
    \end{definition}

    \begin{exemple}
        \begin{enumerate}
            \item Expliquer pourquoi $1$ n'est pas un nombre premier.
            \begin{ZoneReponse}{1cm}
            \end{ZoneReponse}

            \item Parmi les entiers suivants, entourer les nombres premiers :
            \[
                2 \quad 3 \quad 4 \quad 9 \quad 11 \quad 15 \quad 17 \quad 19 \quad 21 \quad 23 \quad 29 \quad 51
            \]
            \begin{ZoneReponse}{2cm}
            \end{ZoneReponse}
        \end{enumerate}
    \end{exemple}

    \begin{methode}[Crible d'Ératosthène]
        Le \textbf{crible d'Ératosthène} est une méthode pour trouver tous les nombres premiers inférieurs à un entier donné $N$ :
        \begin{enumerate}
            \item On écrit tous les entiers de $2$ à $N$ dans un tableau.
            \item On barre le 1 (qui n'est pas premier !).
            \item On entoure $2$ (premier nombre premier), puis on \textbf{barre} tous ses multiples ($4, 6, 8, \ldots$).
            \item On entoure le premier nombre non barré ($3$), puis on barre tous ses multiples non encore barrés.
            \item On répète l'opération jusqu'à ce que tous les entiers soient soit entourés, soit barrés.
            \item Les nombres \textbf{entourés} sont exactement les nombres premiers inférieurs à $N$.
        \end{enumerate}
    \end{methode}

    \begin{exemple}
        Utiliser le crible d'Ératosthène pour trouver tous les nombres premiers inférieurs ou égaux à $100$.

        \medskip
        \begin{center}
        \renewcommand{\arraystretch}{1.8}
        \begin{tabular}{|c|c|c|c|c|c|c|c|c|c|}
        \hline
         1  &  2  &  3  &  4  &  5  &  6  &  7  &  8  &  9  &  10 \\ \hline
        11  & 12  & 13  & 14  & 15  & 16  & 17  & 18  & 19  &  20 \\ \hline
        21  & 22  & 23  & 24  & 25  & 26  & 27  & 28  & 29  &  30 \\ \hline
        31  & 32  & 33  & 34  & 35  & 36  & 37  & 38  & 39  &  40 \\ \hline
        41  & 42  & 43  & 44  & 45  & 46  & 47  & 48  & 49  &  50 \\ \hline
        51  & 52  & 53  & 54  & 55  & 56  & 57  & 58  & 59  &  60 \\ \hline
        61  & 62  & 63  & 64  & 65  & 66  & 67  & 68  & 69  &  70 \\ \hline
        71  & 72  & 73  & 74  & 75  & 76  & 77  & 78  & 79  &  80 \\ \hline
        81  & 82  & 83  & 84  & 85  & 86  & 87  & 88  & 89  &  90 \\ \hline
        91  & 92  & 93  & 94  & 95  & 96  & 97  & 98  & 99  & 100 \\ \hline
        \end{tabular}
        \end{center}

        \medskip
        Conclure : les nombres premiers inférieurs ou égaux à $100$ sont :
        \begin{ZoneReponse}{1.5cm}
        \end{ZoneReponse}
    \end{exemple}


\newpage


% ============================================================
\section{Décomposition d'un entier en produit de facteurs premiers}
% ============================================================

\subsection{Théorème}

    \begin{propriete}[Théorème fondamental de l'arithmétique]
        Tout entier supérieur ou égal à $2$ peut s'écrire de manière \textbf{unique} comme un produit de facteurs premiers (à l'ordre des facteurs près). On appelle cela sa \textbf{décomposition en produit de facteurs premiers}.
    \end{propriete}

    \begin{methode}
        Pour décomposer un entier $n$ en produit de facteurs premiers, on effectue des \textbf{divisions successives} par les nombres premiers ($2, 3, 5, 7, 11, \ldots$) dans l'ordre croissant, jusqu'à obtenir un quotient égal à $1$.
    \end{methode}

    \begin{exemple}
        Décomposer $\np{360}$ en produit de facteurs premiers en utilisant le tableau de divisions successives.

        \medskip
        \begin{minipage}[c]{0.35\textwidth}
            \begin{center}
            \renewcommand{\arraystretch}{1.6}
            \begin{tabular}{r|l}
                \np{360} & \\ \hline
                          & \\ \hline
                          & \\ \hline
                          & \\ \hline
                          & \\ \hline
                          & \\ \hline
                        1 & \\
            \end{tabular}
            \end{center}
        \end{minipage}
        \hfill
        \begin{minipage}[c]{0.58\textwidth}
            \begin{ZoneReponse}{5cm}
            \end{ZoneReponse}
        \end{minipage}
    \end{exemple}

    \begin{exemple}
        Décomposer $\np{2}$~$\np{310}$ en produit de facteurs premiers.

        \medskip
        \begin{minipage}[c]{0.35\textwidth}
            \begin{center}
            \renewcommand{\arraystretch}{1.6}
            \begin{tabular}{r|l}
                \np{2310} & \\ \hline
                           & \\ \hline
                           & \\ \hline
                           & \\ \hline
                           & \\ \hline
                         1 & \\
            \end{tabular}
            \end{center}
        \end{minipage}
        \hfill
        \begin{minipage}[c]{0.58\textwidth}
            \begin{ZoneReponse}{5cm}
            \end{ZoneReponse}
        \end{minipage}
    \end{exemple}


\newpage


\subsection{Application n°1 : rendre une fraction irréductible}

    \begin{definition}
        Une fraction $\dfrac{a}{b}$ est dite \textbf{irréductible} si $a$ et $b$ n'ont pas d'autre diviseur commun que $1$, c'est-à-dire si leur seul diviseur commun est $1$.
    \end{definition}

    \begin{methode}
        Pour rendre une fraction irréductible en utilisant la décomposition en facteurs premiers :
        \begin{enumerate}
            \item On décompose le numérateur et le dénominateur en produit de facteurs premiers.
            \item On repère les facteurs premiers \textbf{communs}.
            \item On simplifie en divisant numérateur et dénominateur par ces facteurs communs.
        \end{enumerate}
    \end{methode}

    \begin{exemple}
        Rendre la fraction $\dfrac{180}{252}$ irréductible.

        \medskip
        \noindent
        \begin{minipage}[t]{0.45\textwidth}
            \textbf{Décomposition de $180$ :}
            \begin{center}
            \renewcommand{\arraystretch}{1.6}
            \begin{tabular}{r|l}
                $180$ & \\ \hline
                      & \\ \hline
                      & \\ \hline
                      & \\ \hline
                    1 & \\
            \end{tabular}
            \end{center}
        \end{minipage}
        \hfill
        \begin{minipage}[t]{0.45\textwidth}
            \textbf{Décomposition de $252$ :}
            \begin{center}
            \renewcommand{\arraystretch}{1.6}
            \begin{tabular}{r|l}
                $252$ & \\ \hline
                      & \\ \hline
                      & \\ \hline
                      & \\ \hline
                    1 & \\
            \end{tabular}
            \end{center}
        \end{minipage}

        \begin{ZoneReponse}{1cm}
        \end{ZoneReponse}
    \end{exemple}

    \begin{exemple}
        Rendre la fraction $\dfrac{\np{1050}}{\np{1260}}$ irréductible.

        \medskip
        \noindent
        \begin{minipage}[t]{0.45\textwidth}
            \textbf{Décomposition de $\np{1050}$ :}
            \begin{center}
            \renewcommand{\arraystretch}{1.6}
            \begin{tabular}{r|l}
                \np{1050} & \\ \hline
                          & \\ \hline
                          & \\ \hline
                          & \\ \hline
                          & \\ \hline
                        1 & \\
            \end{tabular}
            \end{center}
        \end{minipage}
        \hfill
        \begin{minipage}[t]{0.45\textwidth}
            \textbf{Décomposition de $\np{1260}$ :}
            \begin{center}
            \renewcommand{\arraystretch}{1.6}
            \begin{tabular}{r|l}
                \np{1260} & \\ \hline
                          & \\ \hline
                          & \\ \hline
                          & \\ \hline
                          & \\ \hline
                        1 & \\
            \end{tabular}
            \end{center}
        \end{minipage}

        \begin{ZoneReponse}{1cm}
        \end{ZoneReponse}
    \end{exemple}


\newpage


\subsection{Application n°2 : trouver le plus grand diviseur commun à deux entiers}

    \begin{definition}
        Le \textbf{plus grand diviseur commun} (PGCD) de deux entiers $a$ et $b$ est le plus grand entier qui divise à la fois $a$ et $b$. On le note $\text{PGCD}(a\,;\,b)$.
    \end{definition}

    \begin{methode}
        Pour calculer le PGCD de deux entiers à l'aide de leur décomposition en facteurs premiers :
        \begin{enumerate}
            \item On décompose chaque entier en produit de facteurs premiers.
            \item On relève les facteurs premiers \textbf{communs aux deux décompositions}.
            \item Le PGCD est le produit de ces facteurs communs, chacun pris avec le plus \textbf{petit} exposant.
        \end{enumerate}
    \end{methode}

    \begin{exemple}
        Calculer $\text{PGCD}(180\,;\,252)$ en utilisant les décompositions obtenues dans l'exemple précédent.
        \begin{ZoneReponse}{2cm}
        \end{ZoneReponse}
    \end{exemple}

    \begin{remarque}
        Si $\text{PGCD}(a\,;\,b) = 1$, on dit que $a$ et $b$ sont \textbf{premiers entre eux}. Cela signifie qu'ils n'ont aucun facteur premier en commun, et donc que la fraction $\dfrac{a}{b}$ est déjà irréductible.
    \end{remarque}

    \begin{exemple}[Type Brevet – Utilisation du PGCD]
        Une association souhaite constituer des équipes identiques à partir de $\np{168}$ filles et $\np{252}$ garçons. Chaque équipe doit contenir uniquement des filles ou uniquement des garçons, et toutes les équipes doivent avoir le même effectif. On cherche le nombre maximum d'élèves par équipe.

        \begin{enumerate}
            \item Décomposer $\np{168}$ et $\np{252}$ en produit de facteurs premiers.

            \medskip
            \noindent
            \begin{minipage}[t]{0.45\textwidth}
                \textbf{Décomposition de $\np{168}$ :}
                \begin{center}
                \renewcommand{\arraystretch}{1.6}
                \begin{tabular}{r|l}
                    \np{168} & \\ \hline
                             & \\ \hline
                             & \\ \hline
                             & \\ \hline
                           1 & \\
                \end{tabular}
                \end{center}
            \end{minipage}
            \hfill
            \begin{minipage}[t]{0.45\textwidth}
                \textbf{Décomposition de $\np{252}$ :}
                \begin{center}
                \renewcommand{\arraystretch}{1.6}
                \begin{tabular}{r|l}
                    \np{252} & \\ \hline
                             & \\ \hline
                             & \\ \hline
                             & \\ \hline
                           1 & \\
                \end{tabular}
                \end{center}
            \end{minipage}

            \item Calculer $\text{PGCD}(\np{168}\,;\,\np{252})$.
            \begin{ZoneReponse}{2cm}
            \end{ZoneReponse}

            \item En déduire le nombre maximum d'élèves par équipe, ainsi que le nombre d'équipes de filles et d'équipes de garçons.
            \begin{ZoneReponse}{2cm}
            \end{ZoneReponse}
        \end{enumerate}
    \end{exemple}

\end{document}